% !TeX root = ../papers/nc-p3.tex
\section{階数15の分岐曲線全体での中心と長さ六の商}
\label{sec:rank-fifteen-branch-center}
\begin{theorem}\label{thm:rank-fifteen-branch-center}
標数零で、前節の二本の被約分岐曲線の全ての通常点で四次中心の文字次元は
$(3,2,0,0)$ である。文字0と1から二式ずつ選ぶと、四始点で同時に、
長さ0から6の中心商の次元を $1,4,10,20,31,40,44$ にできる。
以下の消去で得られる有限集合は第一表示で2720点、第二表示で2616点あり、
この全5336幾何点は標準Koszul左右AS正則である。
\end{theorem}
\begin{proof}
$F=-t^4+14t^2-1$、$y^2=F$ とする。前節の有理逆を四つの関係空間に代入する。
$\mathbb Q(t,y)$ で分子 $A+Cy$、分母 $q$ を整数係数多項式で持ち、
文字ごとに長さ5まで道を左に延ばす。各消去では元の行の順に独立な行を選び、
ピボットの積でその行とピボット列の小行列式を求める。
この小行列式のノルムの分子、分母、および最終射影の分母の既約因子を記録する。
これで、選んだ小行列式が単元の点では同じ語が基底になる。
全四始点の次数0から5の次元は $1,4,10,20,35,56$ である。
得られた射影で四位置の交換子を計算し、各文字の中心行列にも同じ小行列式の検査を行う。
核の次元は順に $3,2,0,0$ であり、中心の一般基底が特殊化できる範囲が決まる。

両文字から四切断を取り、その係数と元の関係の分母を全て加える。
$t=4,y=13$ の101での還元で短い行から独立なものを選び、
実際の小行列式は標数零の関数体で計算する。
中心商では右に一文字延ばし、二つの連続する始点で長さ6まで
上記の次元を得る。使った小行列式のノルムと分母を再び加える。
これらの行で選んだ語が商を張る。一般正則点での次元と一致するので、
この範囲で特殊化した次元にも等しい。
二位置移動は $(t,y)\mapsto(-1/t,y/t^2)$ である。
実際この変換は $U,r,w,K$ を保ち、$V,n$ の符号を変える。
従って因子集合にこの変換の像を加えれば四始点全てが扱える。
必要なら各始点で有効な中心四式の開集合を同じ選択空間で交叉させる。

通常条件を外れる因子は $t$、$t\pm1$、$t^2+1$、
$t^2\pm2t-1$、$t^2\pm4t+1$、$t^4+10t^2+1$、
$t^8+20t^6-26t^4+20t^2+1$ および無限遠である。
これは $z,\lambda$ と通常条件の各因子の極・零点をノルムで求めて得られる。
$y$ の符号で $w$ は逆数に移り、$\lambda$ は $1/(z\lambda)$ に移るので
通常条件は二枚で一致する。極では $r$ が零または極になり、
無限遠では $K=-1$ となる。
この境界を除いた小行列式の例外集合は次の通りである。
\[
\begin{array}{c|r|r|r}
 \text{表示}&\text{第一座標の既約因子数}&\text{次数の和}&\text{幾何点数}\\ \hline
 1&49&1360&2720\\
 2&47&1308&2616
\end{array}
\]
\path{computations/verify_rank_fifteen_branch_center.py} は因子を関数体から
再生成し、各因子の単純な素体根と、その上の $y$ の相異なる両根を調べる。
各還元で長さ4,5の道、四次中心を求め直し、中心四式の商を全四始点で長さ13まで
計算する。次元は命題\ref{prop:fixed-residual-regularity}の列となる。
各既約因子の単根から得た付値で両根を検査したので、共役および体拡大で全点に戻せる。
同命題の中心持上げと全次数の下界により、表の全点は所要の正則性を持つ。
\end{proof}

\begin{corollary}\label{cor:rank-fifteen-near-collision}
通常倍率平面に $zk^2=1$ の軌跡全体を含む開集合があり、
その上の二つの四重被覆の全窓は標準Koszul左右AS正則である。
\end{corollary}
\begin{proof}
この軌跡上の全窓では定理\ref{thm:rank-fourteen-center}により
道と四次中心が一般次元を持ち、前節までの階数14の結果により中心商に標準次元を
与える四式がある。定階数の近傍に中心四式を延ばせば、長さ13までの
商の次元上界が開近傍で保たれる。道の全次数の下界と合わせて正則性が従う。
各被覆上でこれらの近傍の和を取り、閉補集合の倍率平面への像を除く。
二被覆は有限であり、指定の軌跡の上の全点を近傍が含むので、所要の開集合を得る。
\end{proof}
